\documentclass[11pt,a4paper]{article}
\usepackage[a4paper,margin=1in,headsep=30pt,footskip=35pt]{geometry}
\usepackage[T1]{fontenc}
\usepackage[utf8]{inputenc}
\usepackage{tgpagella}
\usepackage[scale=0.92]{tgheros}
\usepackage{microtype}
\usepackage{graphicx}
\usepackage{hyperref}
\usepackage{enumitem}
\usepackage{xcolor}
\usepackage{titlesec}
\usepackage{parskip}
\usepackage{longtable}
\usepackage{booktabs}
\usepackage{array}
\usepackage{tabularx}
\usepackage{fancyhdr}
\usepackage{lastpage}
\usepackage{tikz}
\usepackage[most]{tcolorbox}

% --- CORE COLOR IDENTITY ---
\definecolor{KazePrimary}{HTML}{284F58}
\definecolor{KazeSecondary}{HTML}{8A6538}
\definecolor{KazeInk}{HTML}{1F292D}
\definecolor{KazeMuted}{HTML}{5D676C}
\definecolor{KazeSoft}{HTML}{F8F4ED}
\definecolor{KazeLine}{HTML}{D8CCBC}
\definecolor{KazePaper}{HTML}{FCFBF8}
\definecolor{KazeLegal}{HTML}{7A2E2E}

\hypersetup{colorlinks=true, linkcolor=KazePrimary, urlcolor=KazePrimary}

\setlist[itemize]{leftmargin=1.5em, itemsep=0.4em}
\renewcommand{\arraystretch}{1.25}
\setlength{\headheight}{35pt}

\newcommand{\KazeLogoPath}{../../composeApp/src/commonMain/composeResources/drawable/k_logo_raster.png}
\newcommand{\KazeMarkPath}{../../composeApp/src/commonMain/composeResources/drawable/k_mark_raster.png}
\newcommand{\GaboLogoPath}{../../composeApp/src/commonMain/composeResources/drawable/gabo_mark_raster.png}
\newcommand{\GaboGrayMarkPath}{../../composeApp/src/commonMain/composeResources/drawable/gabo_mark_gray_raster.png}
\newcommand{\KazeWordmarkHeader}{{\fontfamily{qhv}\selectfont\sffamily\bfseries\small\textcolor{KazePrimary}{KAZE}}}
\newcommand{\KazeWordmarkCover}{{\fontfamily{qhv}\selectfont\sffamily\bfseries\fontsize{30}{30}\selectfont\textcolor{KazePrimary}{KAZE}}}
\newcommand{\KazeByGaboHeader}{{\raisebox{0.05cm}{\includegraphics[height=0.18cm]{\GaboLogoPath}}\hspace{0.2em}\sffamily\tiny\textcolor{KazeSecondary}{by GABO}}}
\newcommand{\KazeByGaboCover}{{\includegraphics[height=0.45cm]{\GaboLogoPath}\hspace{0.4em}\sffamily\large\textcolor{KazeSecondary}{by GABO}}}
\newcommand{\KazeCoverLabel}[1]{{\sffamily\normalsize\textcolor{KazeSecondary}{\textbf{#1}}}}
\newcommand{\KazeCoverTitle}[1]{{\fontfamily{qpl}\selectfont\fontsize{38}{43}\selectfont\textcolor{KazeInk}{#1}}}
\newcommand{\KazeCoverSubtitle}[1]{{\sffamily\Large\textcolor{KazePrimary}{#1}}}
\newcommand{\KazeContactEmail}{orestegabo@icloud.com}
\newcommand{\KazeContactWebsiteUrl}{https://kazerwanda.com}
\newcommand{\KazeContactWebsiteLabel}{kazerwanda.com}
\newcommand{\KazeContactBlock}{
    \begin{itemize}
        \item \textbf{Email:} \href{mailto:\KazeContactEmail}{\KazeContactEmail}
        \item \textbf{Website:} \href{\KazeContactWebsiteUrl}{\KazeContactWebsiteLabel}
    \end{itemize}
}

\pagecolor{KazePaper}

\AddToHook{shipout/background}{%
    \begin{tikzpicture}[remember picture,overlay]
        \fill[KazePrimary, opacity=0.018, rounded corners=14pt]
            ([xshift=0.04\paperwidth,yshift=-0.07\paperheight]current page.north west)
            rectangle
            ([xshift=0.96\paperwidth,yshift=-0.31\paperheight]current page.north west);
        \fill[KazeSecondary, opacity=0.014, rounded corners=14pt]
            ([xshift=0.05\paperwidth,yshift=0.34\paperheight]current page.north west)
            rectangle
            ([xshift=0.95\paperwidth,yshift=0.14\paperheight]current page.south west);

        \draw[KazePrimary, line width=1.7pt, opacity=0.09, line cap=round]
            ([xshift=0.08\paperwidth,yshift=-0.10\paperheight]current page.north west) --
            ([xshift=0.72\paperwidth,yshift=-0.10\paperheight]current page.north west);
        \draw[KazePrimary, line width=0.95pt, opacity=0.085, line cap=round]
            ([xshift=0.12\paperwidth,yshift=-0.135\paperheight]current page.north west) --
            ([xshift=0.88\paperwidth,yshift=-0.135\paperheight]current page.north west);
        \draw[KazeSecondary, line width=1.2pt, opacity=0.08, line cap=round]
            ([xshift=0.18\paperwidth,yshift=0.26\paperheight]current page.south west) --
            ([xshift=0.84\paperwidth,yshift=0.26\paperheight]current page.south west);
        \draw[KazePrimary, line width=1.7pt, opacity=0.09, line cap=round]
            ([xshift=0.14\paperwidth,yshift=0.21\paperheight]current page.south west) --
            ([xshift=0.62\paperwidth,yshift=0.21\paperheight]current page.south west);

        \draw[KazePrimary, line width=1.15pt, opacity=0.09]
            ([xshift=0.62\paperwidth,yshift=-0.06\paperheight]current page.north west) --
            ([xshift=0.76\paperwidth,yshift=-0.06\paperheight]current page.north west) --
            ([xshift=0.82\paperwidth,yshift=-0.11\paperheight]current page.north west) --
            ([xshift=0.93\paperwidth,yshift=-0.11\paperheight]current page.north west) --
            ([xshift=0.93\paperwidth,yshift=-0.18\paperheight]current page.north west) --
            ([xshift=0.82\paperwidth,yshift=-0.18\paperheight]current page.north west) --
            ([xshift=0.75\paperwidth,yshift=-0.24\paperheight]current page.north west) --
            ([xshift=0.58\paperwidth,yshift=-0.24\paperheight]current page.north west);

        \draw[KazeSecondary, line width=1.15pt, opacity=0.08]
            ([xshift=0.10\paperwidth,yshift=0.12\paperheight]current page.south west) --
            ([xshift=0.26\paperwidth,yshift=0.12\paperheight]current page.south west) --
            ([xshift=0.33\paperwidth,yshift=0.17\paperheight]current page.south west) --
            ([xshift=0.48\paperwidth,yshift=0.17\paperheight]current page.south west) --
            ([xshift=0.48\paperwidth,yshift=0.10\paperheight]current page.south west) --
            ([xshift=0.34\paperwidth,yshift=0.10\paperheight]current page.south west) --
            ([xshift=0.27\paperwidth,yshift=0.05\paperheight]current page.south west) --
            ([xshift=0.12\paperwidth,yshift=0.05\paperheight]current page.south west);

        \draw[KazePrimary, line width=1.8pt, opacity=0.05]
            ([xshift=0.70\paperwidth,yshift=-0.22\paperheight]current page.north west) circle [radius=2.0cm];
        \draw[KazeSecondary, line width=1.8pt, opacity=0.048]
            ([xshift=0.14\paperwidth,yshift=0.18\paperheight]current page.south west) circle [radius=2.35cm];

        \draw[KazePrimary, line cap=round, opacity=0.065, line width=1.2pt]
            ([xshift=0.11\paperwidth,yshift=-0.185\paperheight]current page.north west) --
            ([xshift=0.36\paperwidth,yshift=-0.185\paperheight]current page.north west);
        \draw[KazeSecondary, line cap=round, opacity=0.06, line width=0.8pt]
            ([xshift=0.11\paperwidth,yshift=-0.203\paperheight]current page.north west) --
            ([xshift=0.42\paperwidth,yshift=-0.203\paperheight]current page.north west);
        \draw[KazePrimary, line cap=round, opacity=0.065, line width=1.2pt]
            ([xshift=0.11\paperwidth,yshift=-0.221\paperheight]current page.north west) --
            ([xshift=0.48\paperwidth,yshift=-0.221\paperheight]current page.north west);
        \draw[KazeSecondary, line cap=round, opacity=0.06, line width=0.8pt]
            ([xshift=0.11\paperwidth,yshift=-0.239\paperheight]current page.north west) --
            ([xshift=0.54\paperwidth,yshift=-0.239\paperheight]current page.north west);
        \draw[KazePrimary, line cap=round, opacity=0.065, line width=1.2pt]
            ([xshift=0.11\paperwidth,yshift=-0.257\paperheight]current page.north west) --
            ([xshift=0.60\paperwidth,yshift=-0.257\paperheight]current page.north west);
        \draw[KazeSecondary, line cap=round, opacity=0.06, line width=0.8pt]
            ([xshift=0.11\paperwidth,yshift=-0.275\paperheight]current page.north west) --
            ([xshift=0.66\paperwidth,yshift=-0.275\paperheight]current page.north west);
        \draw[KazePrimary, line cap=round, opacity=0.065, line width=1.2pt]
            ([xshift=0.11\paperwidth,yshift=-0.293\paperheight]current page.north west) --
            ([xshift=0.72\paperwidth,yshift=-0.293\paperheight]current page.north west);
        \node[anchor=south east, opacity=0.035] at ([xshift=-0.45cm,yshift=0.65cm]current page.south east) {%
            \includegraphics[width=4.8cm]{\KazeMarkPath}%
        };
        \node[anchor=north east, opacity=0.028] at ([xshift=-1.1cm,yshift=-1.2cm]current page.north east) {%
            \includegraphics[width=2.2cm]{\GaboGrayMarkPath}%
        };
    \end{tikzpicture}%
}

% --- SECTION DESIGN: TECHNICAL PRECISION ---
\titleformat{\section}
{\sffamily\Large\bfseries\color{KazeInk}}
{\thesection}
{1em}
{\MakeUppercase}
[\vspace{0.2em}{\color{KazeSecondary}\hrule height 1.5pt width 2cm}\vspace{0.5em}]

\titleformat{\subsection}
{\sffamily\large\bfseries\color{KazePrimary}}
{\thesubsection}
{1em}
{}

% --- DUAL-BRANDED HEADER (HIGH VISIBILITY) ---
\pagestyle{fancy}
\fancyhf{}
\renewcommand{\headrulewidth}{0pt}
\renewcommand{\footrulewidth}{0pt}

\fancyhead[L]{%
    \begin{tabular}{@{}l l@{}}
        \raisebox{-0.3\height}{\includegraphics[height=0.6cm]{\KazeLogoPath}} &
        \begin{minipage}[b]{4cm}
            \KazeWordmarkHeader\\
            \vspace{-1mm}
            \KazeByGaboHeader
        \end{minipage}
    \end{tabular}
}
\fancyhead[R]{\sffamily\tiny\textcolor{KazeMuted}{KAZE PROPRIETARY // \today}}

\fancyfoot[L]{%
    \begin{tabular}{@{}l@{}}
        {\color{KazeLine}\rule{3.1cm}{0.6pt}}\\[-0.2em]
        {\sffamily\tiny\textcolor{KazeLegal}{CONFIDENTIAL PRIVATE DOCUMENT}}
    \end{tabular}
}
\fancyfoot[C]{%
    \sffamily\scriptsize
    \textcolor{KazeMuted}{\textsc{Page} \thepage\ of \pageref*{LastPage}}
}
\fancyfoot[R]{%
    \begin{tabular}{@{}r@{}}
        {\color{KazeLine}\rule{2.1cm}{0.6pt}}\\[-0.2em]
        {\sffamily\tiny\textcolor{KazePrimary}{Kaze by GABO}}
    \end{tabular}
}

% --- CLEAN REGISTRY BOXES ---
\newtcolorbox{kazehighlightbox}{
    enhanced,
    colback=KazeSoft,      % Only a very faint paper-like tint
    frame hidden,          % No borders at all
    sharp corners,
    left=5mm, right=5mm, top=4mm, bottom=4mm,
% Use a small square bullet as an anchor instead of a line
    overlay={
        \fill[KazeSecondary] ([xshift=2mm,yshift=-4mm]frame.north west) rectangle ([xshift=3mm,yshift=-5mm]frame.north west);
    }
}
\newtcolorbox{kazecalloutbox}[1][]{
    enhanced,
    colback=white,
    colframe=KazeInk,
    boxrule=1pt,
    sharp corners,
    fonttitle=\sffamily\bfseries\small,
    coltitle=white,
    attach boxed title to top left={yshift=-2mm, xshift=5mm},
    boxed title style={colback=KazeInk, sharp corners, boxrule=0pt},
    title=#1,
    left=5mm, top=6mm, bottom=4mm
}

\newcommand{\KazePageTitle}[1]{%
    \vspace{1em}
    {\sffamily\small\textcolor{KazeSecondary}{\textbf{\MakeUppercase{Document Scope: #1}}}\par}
    \vspace{0.5em}
}

\newcommand{\KazeDocumentNotice}{%
    \begin{kazecalloutbox}[Security Disclosure]
    \sffamily\small This dossier contains trade secrets and sensitive intellectual property. \textbf{Unauthorized distribution is strictly prohibited.} Access is limited to Kaze personnel and authorized external stakeholders for legal and investor review.
    \end{kazecalloutbox}
}

% --- THE "Blueprints" COVER PAGE ---
\newcommand{\KazeCover}[4]{
    \hypersetup{pageanchor=false}
    \begin{titlepage}
    \thispagestyle{empty}

% Background "Pillar" for institutional weight
    \begin{tikzpicture}[remember picture,overlay]
    \fill[KazePrimary] (current page.north west) rectangle ([xshift=0.5cm]current page.south west);
    \draw[KazePrimary!55, line width=1.8pt, opacity=0.22, line cap=round]
        ([xshift=0.08\paperwidth,yshift=-0.10\paperheight]current page.north west) --
        ([xshift=0.72\paperwidth,yshift=-0.10\paperheight]current page.north west);
    \draw[KazeSecondary!70, line width=1.1pt, opacity=0.19, line cap=round]
        ([xshift=0.12\paperwidth,yshift=-0.135\paperheight]current page.north west) --
        ([xshift=0.88\paperwidth,yshift=-0.135\paperheight]current page.north west);
    \draw[KazePrimary!55, line width=1.2pt, opacity=0.20]
        ([xshift=0.62\paperwidth,yshift=-0.06\paperheight]current page.north west) --
        ([xshift=0.76\paperwidth,yshift=-0.06\paperheight]current page.north west) --
        ([xshift=0.82\paperwidth,yshift=-0.11\paperheight]current page.north west) --
        ([xshift=0.93\paperwidth,yshift=-0.11\paperheight]current page.north west) --
        ([xshift=0.93\paperwidth,yshift=-0.18\paperheight]current page.north west) --
        ([xshift=0.82\paperwidth,yshift=-0.18\paperheight]current page.north west) --
        ([xshift=0.75\paperwidth,yshift=-0.24\paperheight]current page.north west) --
        ([xshift=0.58\paperwidth,yshift=-0.24\paperheight]current page.north west);
    \draw[KazeSecondary!75, line width=1.35pt, opacity=0.18, line cap=round]
        ([xshift=0.18\paperwidth,yshift=0.26\paperheight]current page.south west) --
        ([xshift=0.84\paperwidth,yshift=0.26\paperheight]current page.south west);
    \draw[KazePrimary!55, line width=1.8pt, opacity=0.20, line cap=round]
        ([xshift=0.14\paperwidth,yshift=0.21\paperheight]current page.south west) --
        ([xshift=0.62\paperwidth,yshift=0.21\paperheight]current page.south west);
    \end{tikzpicture}

    \vspace*{0.5cm}
    \begin{flushleft}
    \hspace{1cm}
    \begin{tabular}{@{}l l@{}}
    \raisebox{-0.3\height}{\includegraphics[width=2.2cm]{\KazeLogoPath}} &
    \begin{minipage}[b]{8cm}
    {\KazeWordmarkCover\par}
    \vspace{0.1cm}
    {\KazeByGaboCover\par}
    \end{minipage}
    \end{tabular}
    \end{flushleft}

    \vspace{3.5cm}

    \hspace{1cm}
    \begin{minipage}{0.85\textwidth}
    \raggedright
    {\KazeCoverLabel{OFFICIAL DOSSIER / #4}\par}
    \vspace{0.95cm}
    {\KazeCoverTitle{#1}\par}
    \vspace{0.7cm}
    {\KazeCoverSubtitle{#2}\par}

    \vspace{1.5cm}
    {\color{KazeLine}\hrule height 2pt width 3cm}
    \vspace{1.5cm}

    {\sffamily\large\textcolor{KazeMuted}{#3}\par}
    \end{minipage}

    \vfill

    \hspace{1cm}
    \begin{minipage}{0.85\textwidth}
    \begin{kazehighlightbox}
    \sffamily\small
    \textbf{\textcolor{KazeLegal}{RESTRICTED DOCUMENT:}} This information is intended for the recipient only. It contains proprietary strategic planning belonging to Kaze. By accessing this document, you agree to non-disclosure terms.
    \end{kazehighlightbox}
    \end{minipage}

    \vspace{1.5cm}
    \hspace{1cm}
    \begin{minipage}{0.85\textwidth}
    \sffamily\tiny\textcolor{KazeMuted}{AUTHORIZATION \& CONTACT:} \\
    \sffamily\small\textbf{\href{mailto:\KazeContactEmail}{\KazeContactEmail}} \hfill \textbf{\MakeUppercase{\KazeContactWebsiteLabel}}
    \end{minipage}

    \end{titlepage}
    \hypersetup{pageanchor=true}
}

\newtcblisting{kazecodeblock}{
  enhanced,
  breakable,
  colback=white,
  colframe=KazeLine,
  boxrule=0.6pt,
  sharp corners,
  left=3mm,
  right=3mm,
  top=2mm,
  bottom=2mm,
  listing only,
  listing options={
    basicstyle=\ttfamily\small,
    breaklines=true,
    columns=fullflexible
  }
}
\newtcolorbox{kazeexamplebox}[1][]{
  enhanced,
  breakable,
  colback=white,
  colframe=KazeLine,
  boxrule=0.6pt,
  sharp corners,
  fonttitle=\sffamily\bfseries\small,
  coltitle=KazeInk,
  title=#1,
  left=4mm,
  right=4mm,
  top=3mm,
  bottom=3mm
}

\begin{document}

\KazeCover
  {Kaze PDF Cheatsheet}
  {Copy-paste guide for creating themed LaTeX documents}
  {April 2026}
  {Internal documentation aid}

\KazePageTitle{PDF Cheatsheet}
\KazeDocumentNotice

\section{What This PDF Is For}

This document shows the current Kaze LaTeX theme piece by piece so a new document can be assembled quickly.

\begin{kazecalloutbox}[Recommended Use]
Keep this PDF open while writing a new LaTeX document. Copy the relevant block, paste it into a new file, then replace the placeholder content.
\end{kazecalloutbox}

\section{Smallest Working File}

Start from this exact structure:

\begin{kazecodeblock}
\documentclass[11pt,a4paper]{article}
\usepackage[a4paper,margin=1in,headsep=30pt,footskip=35pt]{geometry}
\usepackage[T1]{fontenc}
\usepackage[utf8]{inputenc}
\usepackage{tgpagella}
\usepackage[scale=0.92]{tgheros}
\usepackage{microtype}
\usepackage{graphicx}
\usepackage{hyperref}
\usepackage{enumitem}
\usepackage{xcolor}
\usepackage{titlesec}
\usepackage{parskip}
\usepackage{longtable}
\usepackage{booktabs}
\usepackage{array}
\usepackage{tabularx}
\usepackage{fancyhdr}
\usepackage{lastpage}
\usepackage{tikz}
\usepackage[most]{tcolorbox}

% --- CORE COLOR IDENTITY ---
\definecolor{KazePrimary}{HTML}{284F58}
\definecolor{KazeSecondary}{HTML}{8A6538}
\definecolor{KazeInk}{HTML}{1F292D}
\definecolor{KazeMuted}{HTML}{5D676C}
\definecolor{KazeSoft}{HTML}{F8F4ED}
\definecolor{KazeLine}{HTML}{D8CCBC}
\definecolor{KazePaper}{HTML}{FCFBF8}
\definecolor{KazeLegal}{HTML}{7A2E2E}

\hypersetup{colorlinks=true, linkcolor=KazePrimary, urlcolor=KazePrimary}

\setlist[itemize]{leftmargin=1.5em, itemsep=0.4em}
\renewcommand{\arraystretch}{1.25}
\setlength{\headheight}{35pt}

\newcommand{\KazeLogoPath}{../../composeApp/src/commonMain/composeResources/drawable/k_logo_raster.png}
\newcommand{\KazeMarkPath}{../../composeApp/src/commonMain/composeResources/drawable/k_mark_raster.png}
\newcommand{\GaboLogoPath}{../../composeApp/src/commonMain/composeResources/drawable/gabo_mark_raster.png}
\newcommand{\GaboGrayMarkPath}{../../composeApp/src/commonMain/composeResources/drawable/gabo_mark_gray_raster.png}
\newcommand{\KazeWordmarkHeader}{{\fontfamily{qhv}\selectfont\sffamily\bfseries\small\textcolor{KazePrimary}{KAZE}}}
\newcommand{\KazeWordmarkCover}{{\fontfamily{qhv}\selectfont\sffamily\bfseries\fontsize{30}{30}\selectfont\textcolor{KazePrimary}{KAZE}}}
\newcommand{\KazeByGaboHeader}{{\raisebox{0.05cm}{\includegraphics[height=0.18cm]{\GaboLogoPath}}\hspace{0.2em}\sffamily\tiny\textcolor{KazeSecondary}{by GABO}}}
\newcommand{\KazeByGaboCover}{{\includegraphics[height=0.45cm]{\GaboLogoPath}\hspace{0.4em}\sffamily\large\textcolor{KazeSecondary}{by GABO}}}
\newcommand{\KazeCoverLabel}[1]{{\sffamily\normalsize\textcolor{KazeSecondary}{\textbf{#1}}}}
\newcommand{\KazeCoverTitle}[1]{{\fontfamily{qpl}\selectfont\fontsize{38}{43}\selectfont\textcolor{KazeInk}{#1}}}
\newcommand{\KazeCoverSubtitle}[1]{{\sffamily\Large\textcolor{KazePrimary}{#1}}}
\newcommand{\KazeContactEmail}{orestegabo@icloud.com}
\newcommand{\KazeContactWebsiteUrl}{https://kazerwanda.com}
\newcommand{\KazeContactWebsiteLabel}{kazerwanda.com}
\newcommand{\KazeContactBlock}{
    \begin{itemize}
        \item \textbf{Email:} \href{mailto:\KazeContactEmail}{\KazeContactEmail}
        \item \textbf{Website:} \href{\KazeContactWebsiteUrl}{\KazeContactWebsiteLabel}
    \end{itemize}
}

\pagecolor{KazePaper}

\AddToHook{shipout/background}{%
    \begin{tikzpicture}[remember picture,overlay]
        \fill[KazePrimary, opacity=0.018, rounded corners=14pt]
            ([xshift=0.04\paperwidth,yshift=-0.07\paperheight]current page.north west)
            rectangle
            ([xshift=0.96\paperwidth,yshift=-0.31\paperheight]current page.north west);
        \fill[KazeSecondary, opacity=0.014, rounded corners=14pt]
            ([xshift=0.05\paperwidth,yshift=0.34\paperheight]current page.north west)
            rectangle
            ([xshift=0.95\paperwidth,yshift=0.14\paperheight]current page.south west);

        \draw[KazePrimary, line width=1.7pt, opacity=0.09, line cap=round]
            ([xshift=0.08\paperwidth,yshift=-0.10\paperheight]current page.north west) --
            ([xshift=0.72\paperwidth,yshift=-0.10\paperheight]current page.north west);
        \draw[KazePrimary, line width=0.95pt, opacity=0.085, line cap=round]
            ([xshift=0.12\paperwidth,yshift=-0.135\paperheight]current page.north west) --
            ([xshift=0.88\paperwidth,yshift=-0.135\paperheight]current page.north west);
        \draw[KazeSecondary, line width=1.2pt, opacity=0.08, line cap=round]
            ([xshift=0.18\paperwidth,yshift=0.26\paperheight]current page.south west) --
            ([xshift=0.84\paperwidth,yshift=0.26\paperheight]current page.south west);
        \draw[KazePrimary, line width=1.7pt, opacity=0.09, line cap=round]
            ([xshift=0.14\paperwidth,yshift=0.21\paperheight]current page.south west) --
            ([xshift=0.62\paperwidth,yshift=0.21\paperheight]current page.south west);

        \draw[KazePrimary, line width=1.15pt, opacity=0.09]
            ([xshift=0.62\paperwidth,yshift=-0.06\paperheight]current page.north west) --
            ([xshift=0.76\paperwidth,yshift=-0.06\paperheight]current page.north west) --
            ([xshift=0.82\paperwidth,yshift=-0.11\paperheight]current page.north west) --
            ([xshift=0.93\paperwidth,yshift=-0.11\paperheight]current page.north west) --
            ([xshift=0.93\paperwidth,yshift=-0.18\paperheight]current page.north west) --
            ([xshift=0.82\paperwidth,yshift=-0.18\paperheight]current page.north west) --
            ([xshift=0.75\paperwidth,yshift=-0.24\paperheight]current page.north west) --
            ([xshift=0.58\paperwidth,yshift=-0.24\paperheight]current page.north west);

        \draw[KazeSecondary, line width=1.15pt, opacity=0.08]
            ([xshift=0.10\paperwidth,yshift=0.12\paperheight]current page.south west) --
            ([xshift=0.26\paperwidth,yshift=0.12\paperheight]current page.south west) --
            ([xshift=0.33\paperwidth,yshift=0.17\paperheight]current page.south west) --
            ([xshift=0.48\paperwidth,yshift=0.17\paperheight]current page.south west) --
            ([xshift=0.48\paperwidth,yshift=0.10\paperheight]current page.south west) --
            ([xshift=0.34\paperwidth,yshift=0.10\paperheight]current page.south west) --
            ([xshift=0.27\paperwidth,yshift=0.05\paperheight]current page.south west) --
            ([xshift=0.12\paperwidth,yshift=0.05\paperheight]current page.south west);

        \draw[KazePrimary, line width=1.8pt, opacity=0.05]
            ([xshift=0.70\paperwidth,yshift=-0.22\paperheight]current page.north west) circle [radius=2.0cm];
        \draw[KazeSecondary, line width=1.8pt, opacity=0.048]
            ([xshift=0.14\paperwidth,yshift=0.18\paperheight]current page.south west) circle [radius=2.35cm];

        \draw[KazePrimary, line cap=round, opacity=0.065, line width=1.2pt]
            ([xshift=0.11\paperwidth,yshift=-0.185\paperheight]current page.north west) --
            ([xshift=0.36\paperwidth,yshift=-0.185\paperheight]current page.north west);
        \draw[KazeSecondary, line cap=round, opacity=0.06, line width=0.8pt]
            ([xshift=0.11\paperwidth,yshift=-0.203\paperheight]current page.north west) --
            ([xshift=0.42\paperwidth,yshift=-0.203\paperheight]current page.north west);
        \draw[KazePrimary, line cap=round, opacity=0.065, line width=1.2pt]
            ([xshift=0.11\paperwidth,yshift=-0.221\paperheight]current page.north west) --
            ([xshift=0.48\paperwidth,yshift=-0.221\paperheight]current page.north west);
        \draw[KazeSecondary, line cap=round, opacity=0.06, line width=0.8pt]
            ([xshift=0.11\paperwidth,yshift=-0.239\paperheight]current page.north west) --
            ([xshift=0.54\paperwidth,yshift=-0.239\paperheight]current page.north west);
        \draw[KazePrimary, line cap=round, opacity=0.065, line width=1.2pt]
            ([xshift=0.11\paperwidth,yshift=-0.257\paperheight]current page.north west) --
            ([xshift=0.60\paperwidth,yshift=-0.257\paperheight]current page.north west);
        \draw[KazeSecondary, line cap=round, opacity=0.06, line width=0.8pt]
            ([xshift=0.11\paperwidth,yshift=-0.275\paperheight]current page.north west) --
            ([xshift=0.66\paperwidth,yshift=-0.275\paperheight]current page.north west);
        \draw[KazePrimary, line cap=round, opacity=0.065, line width=1.2pt]
            ([xshift=0.11\paperwidth,yshift=-0.293\paperheight]current page.north west) --
            ([xshift=0.72\paperwidth,yshift=-0.293\paperheight]current page.north west);
        \node[anchor=south east, opacity=0.035] at ([xshift=-0.45cm,yshift=0.65cm]current page.south east) {%
            \includegraphics[width=4.8cm]{\KazeMarkPath}%
        };
        \node[anchor=north east, opacity=0.028] at ([xshift=-1.1cm,yshift=-1.2cm]current page.north east) {%
            \includegraphics[width=2.2cm]{\GaboGrayMarkPath}%
        };
    \end{tikzpicture}%
}

% --- SECTION DESIGN: TECHNICAL PRECISION ---
\titleformat{\section}
{\sffamily\Large\bfseries\color{KazeInk}}
{\thesection}
{1em}
{\MakeUppercase}
[\vspace{0.2em}{\color{KazeSecondary}\hrule height 1.5pt width 2cm}\vspace{0.5em}]

\titleformat{\subsection}
{\sffamily\large\bfseries\color{KazePrimary}}
{\thesubsection}
{1em}
{}

% --- DUAL-BRANDED HEADER (HIGH VISIBILITY) ---
\pagestyle{fancy}
\fancyhf{}
\renewcommand{\headrulewidth}{0pt}
\renewcommand{\footrulewidth}{0pt}

\fancyhead[L]{%
    \begin{tabular}{@{}l l@{}}
        \raisebox{-0.3\height}{\includegraphics[height=0.6cm]{\KazeLogoPath}} &
        \begin{minipage}[b]{4cm}
            \KazeWordmarkHeader\\
            \vspace{-1mm}
            \KazeByGaboHeader
        \end{minipage}
    \end{tabular}
}
\fancyhead[R]{\sffamily\tiny\textcolor{KazeMuted}{KAZE PROPRIETARY // \today}}

\fancyfoot[L]{%
    \begin{tabular}{@{}l@{}}
        {\color{KazeLine}\rule{3.1cm}{0.6pt}}\\[-0.2em]
        {\sffamily\tiny\textcolor{KazeLegal}{CONFIDENTIAL PRIVATE DOCUMENT}}
    \end{tabular}
}
\fancyfoot[C]{%
    \sffamily\scriptsize
    \textcolor{KazeMuted}{\textsc{Page} \thepage\ of \pageref*{LastPage}}
}
\fancyfoot[R]{%
    \begin{tabular}{@{}r@{}}
        {\color{KazeLine}\rule{2.1cm}{0.6pt}}\\[-0.2em]
        {\sffamily\tiny\textcolor{KazePrimary}{Kaze by GABO}}
    \end{tabular}
}

% --- CLEAN REGISTRY BOXES ---
\newtcolorbox{kazehighlightbox}{
    enhanced,
    colback=KazeSoft,      % Only a very faint paper-like tint
    frame hidden,          % No borders at all
    sharp corners,
    left=5mm, right=5mm, top=4mm, bottom=4mm,
% Use a small square bullet as an anchor instead of a line
    overlay={
        \fill[KazeSecondary] ([xshift=2mm,yshift=-4mm]frame.north west) rectangle ([xshift=3mm,yshift=-5mm]frame.north west);
    }
}
\newtcolorbox{kazecalloutbox}[1][]{
    enhanced,
    colback=white,
    colframe=KazeInk,
    boxrule=1pt,
    sharp corners,
    fonttitle=\sffamily\bfseries\small,
    coltitle=white,
    attach boxed title to top left={yshift=-2mm, xshift=5mm},
    boxed title style={colback=KazeInk, sharp corners, boxrule=0pt},
    title=#1,
    left=5mm, top=6mm, bottom=4mm
}

\newcommand{\KazePageTitle}[1]{%
    \vspace{1em}
    {\sffamily\small\textcolor{KazeSecondary}{\textbf{\MakeUppercase{Document Scope: #1}}}\par}
    \vspace{0.5em}
}

\newcommand{\KazeDocumentNotice}{%
    \begin{kazecalloutbox}[Security Disclosure]
    \sffamily\small This dossier contains trade secrets and sensitive intellectual property. \textbf{Unauthorized distribution is strictly prohibited.} Access is limited to Kaze personnel and authorized external stakeholders for legal and investor review.
    \end{kazecalloutbox}
}

% --- THE "Blueprints" COVER PAGE ---
\newcommand{\KazeCover}[4]{
    \hypersetup{pageanchor=false}
    \begin{titlepage}
    \thispagestyle{empty}

% Background "Pillar" for institutional weight
    \begin{tikzpicture}[remember picture,overlay]
    \fill[KazePrimary] (current page.north west) rectangle ([xshift=0.5cm]current page.south west);
    \draw[KazePrimary!55, line width=1.8pt, opacity=0.22, line cap=round]
        ([xshift=0.08\paperwidth,yshift=-0.10\paperheight]current page.north west) --
        ([xshift=0.72\paperwidth,yshift=-0.10\paperheight]current page.north west);
    \draw[KazeSecondary!70, line width=1.1pt, opacity=0.19, line cap=round]
        ([xshift=0.12\paperwidth,yshift=-0.135\paperheight]current page.north west) --
        ([xshift=0.88\paperwidth,yshift=-0.135\paperheight]current page.north west);
    \draw[KazePrimary!55, line width=1.2pt, opacity=0.20]
        ([xshift=0.62\paperwidth,yshift=-0.06\paperheight]current page.north west) --
        ([xshift=0.76\paperwidth,yshift=-0.06\paperheight]current page.north west) --
        ([xshift=0.82\paperwidth,yshift=-0.11\paperheight]current page.north west) --
        ([xshift=0.93\paperwidth,yshift=-0.11\paperheight]current page.north west) --
        ([xshift=0.93\paperwidth,yshift=-0.18\paperheight]current page.north west) --
        ([xshift=0.82\paperwidth,yshift=-0.18\paperheight]current page.north west) --
        ([xshift=0.75\paperwidth,yshift=-0.24\paperheight]current page.north west) --
        ([xshift=0.58\paperwidth,yshift=-0.24\paperheight]current page.north west);
    \draw[KazeSecondary!75, line width=1.35pt, opacity=0.18, line cap=round]
        ([xshift=0.18\paperwidth,yshift=0.26\paperheight]current page.south west) --
        ([xshift=0.84\paperwidth,yshift=0.26\paperheight]current page.south west);
    \draw[KazePrimary!55, line width=1.8pt, opacity=0.20, line cap=round]
        ([xshift=0.14\paperwidth,yshift=0.21\paperheight]current page.south west) --
        ([xshift=0.62\paperwidth,yshift=0.21\paperheight]current page.south west);
    \end{tikzpicture}

    \vspace*{0.5cm}
    \begin{flushleft}
    \hspace{1cm}
    \begin{tabular}{@{}l l@{}}
    \raisebox{-0.3\height}{\includegraphics[width=2.2cm]{\KazeLogoPath}} &
    \begin{minipage}[b]{8cm}
    {\KazeWordmarkCover\par}
    \vspace{0.1cm}
    {\KazeByGaboCover\par}
    \end{minipage}
    \end{tabular}
    \end{flushleft}

    \vspace{3.5cm}

    \hspace{1cm}
    \begin{minipage}{0.85\textwidth}
    \raggedright
    {\KazeCoverLabel{OFFICIAL DOSSIER / #4}\par}
    \vspace{0.95cm}
    {\KazeCoverTitle{#1}\par}
    \vspace{0.7cm}
    {\KazeCoverSubtitle{#2}\par}

    \vspace{1.5cm}
    {\color{KazeLine}\hrule height 2pt width 3cm}
    \vspace{1.5cm}

    {\sffamily\large\textcolor{KazeMuted}{#3}\par}
    \end{minipage}

    \vfill

    \hspace{1cm}
    \begin{minipage}{0.85\textwidth}
    \begin{kazehighlightbox}
    \sffamily\small
    \textbf{\textcolor{KazeLegal}{RESTRICTED DOCUMENT:}} This information is intended for the recipient only. It contains proprietary strategic planning belonging to Kaze. By accessing this document, you agree to non-disclosure terms.
    \end{kazehighlightbox}
    \end{minipage}

    \vspace{1.5cm}
    \hspace{1cm}
    \begin{minipage}{0.85\textwidth}
    \sffamily\tiny\textcolor{KazeMuted}{AUTHORIZATION \& CONTACT:} \\
    \sffamily\small\textbf{\href{mailto:\KazeContactEmail}{\KazeContactEmail}} \hfill \textbf{\MakeUppercase{\KazeContactWebsiteLabel}}
    \end{minipage}

    \end{titlepage}
    \hypersetup{pageanchor=true}
}


\begin{document}

\KazeCover
  {Document Title}
  {Short descriptive subtitle}
  {April 2026}
  {Audience or document type}

\KazePageTitle{Document Scope}
\KazeDocumentNotice

\section{First Section}

Your content here.

\end{document}
\end{kazecodeblock}

\begin{kazehighlightbox}
\textbf{Use this first:} If you only need one snippet from this cheatsheet, use the block above. It is the smallest complete file that matches the current Kaze theme.
\end{kazehighlightbox}

\section{Cover Block}

Use this near the top of the document:

\begin{kazecodeblock}
\KazeCover
  {Kaze Investor Brief}
  {Venue Experience, Access, and Commerce Platform}
  {April 2026}
  {Investor and partner-facing summary}
\end{kazecodeblock}

Arguments:
\begin{itemize}
    \item first line: main title
    \item second line: subtitle
    \item third line: date
    \item fourth line: document type or intended audience
\end{itemize}

\begin{kazeexamplebox}[Result Preview]
\textbf{Visual effect:} a branded first page with the Kaze by GABO lockup, the main document title, subtitle, date, private-use notice, and contact line.
\end{kazeexamplebox}

\section{Document Scope and Notice}

Use these directly after the cover:

\begin{kazecodeblock}
\KazePageTitle{Investor Brief}
\KazeDocumentNotice
\end{kazecodeblock}

This gives the body a scope label and the standard confidentiality box.

\begin{kazeexamplebox}[Result Preview]
\KazePageTitle{Investor Brief}
\KazeDocumentNotice
\end{kazeexamplebox}

\section{Normal Sections}

\subsection{Section Heading}

\begin{kazecodeblock}
\section{Executive Summary}

Kaze is a venue experience, access, and commerce platform.
\end{kazecodeblock}

\begin{kazeexamplebox}[Result Preview]
\section*{Executive Summary}
Kaze is a venue experience, access, and commerce platform.
\end{kazeexamplebox}

\subsection{Subsection Heading}

\begin{kazecodeblock}
\subsection{Payments}

Kaze can support MTN MoMo, Airtel Money, and other local rails.
\end{kazecodeblock}

\begin{kazeexamplebox}[Result Preview]
\subsection*{Payments}
Kaze can support MTN MoMo, Airtel Money, and other local rails.
\end{kazeexamplebox}

\section{Highlight Box}

Use this for strategic notes or key observations:

\begin{kazecodeblock}
\begin{kazehighlightbox}
\textbf{Strategic note:} Kaze can reuse venue-map infrastructure
across hotels, conferences, and weddings.
\end{kazehighlightbox}
\end{kazecodeblock}

\begin{kazeexamplebox}[Result Preview]
\begin{kazehighlightbox}
\textbf{Strategic note:} Kaze can reuse venue-map infrastructure across hotels, conferences, and weddings.
\end{kazehighlightbox}
\end{kazeexamplebox}

\section{Callout Box With Title}

Use this for summary panels and explanatory notes:

\begin{kazecodeblock}
\begin{kazecalloutbox}[Why This Matters]
This document helps future developers understand the architecture
before reading the codebase.
\end{kazecalloutbox}
\end{kazecodeblock}

\begin{kazeexamplebox}[Result Preview]
\begin{kazecalloutbox}[Why This Matters]
This document helps future developers understand the architecture before reading the codebase.
\end{kazecalloutbox}
\end{kazeexamplebox}

\section{Lists}

\begin{kazecodeblock}
\begin{itemize}
    \item hotels
    \item conference venues
    \item wedding venues
\end{itemize}
\end{kazecodeblock}

\begin{kazeexamplebox}[Result Preview]
\begin{itemize}
    \item hotels
    \item conference venues
    \item wedding venues
\end{itemize}
\end{kazeexamplebox}

\section{Simple Table Layout}

Use this for structured comparisons or opportunity areas:

\begin{kazecodeblock}
\begin{longtable}{p{0.28\textwidth} p{0.64\textwidth}}
\textbf{Area} & \textbf{Why It Matters} \\
Hotels & strong starting market for guest flows and services \\
Conference venues & good fit for bookings and access control \\
\end{longtable}
\end{kazecodeblock}

\begin{kazeexamplebox}[Result Preview]
\begin{longtable}{p{0.28\textwidth} p{0.56\textwidth}}
\textbf{Area} & \textbf{Why It Matters} \\
Hotels & strong starting market for guest flows and services \\
Conference venues & good fit for bookings and access control \\
\end{longtable}
\end{kazeexamplebox}

\section{Contact Block}

\begin{kazecodeblock}
\section{Contact}

\begin{itemize}
    \item \textbf{Email:} \href{mailto:orestegabo@icloud.com}{orestegabo@icloud.com}
    \item \textbf{Website:} \href{https://orestegabo.dev}{https://orestegabo.dev}
\end{itemize}
\end{kazecodeblock}

\begin{kazeexamplebox}[Result Preview]
\begin{itemize}
    \item \textbf{Email:} \href{mailto:orestegabo@icloud.com}{orestegabo@icloud.com}
    \item \textbf{Website:} \href{https://orestegabo.dev}{orestegabo.dev}
\end{itemize}
\end{kazeexamplebox}

\section{Recommended Workflow}

\begin{itemize}
    \item duplicate `kaze-template.tex` or start from the minimal file in this document
    \item rename the file clearly
    \item replace the cover values
    \item write sections
    \item add highlight and callout boxes where needed
    \item build with `sh build-docs.sh`
\end{itemize}

\section{File Naming Advice}

Good file names:
\begin{itemize}
    \item `kaze-investor-memo.tex`
    \item `kaze-business-plan.tex`
    \item `kaze-architecture-spec.tex`
    \item `kaze-partner-brief.tex`
\end{itemize}

Avoid:
\begin{itemize}
    \item `doc1.tex`
    \item `notes.tex`
    \item `final-final.tex`
\end{itemize}

\section{Where To Look Next}

For source references:
\begin{itemize}
    \item see `CHEATSHEET.md` for the Markdown version
    \item see `kaze-template.tex` for a reusable starter file
    \item see `kaze-theme.tex` only if you want to modify the global design
\end{itemize}

\end{document}
